\section{Case Study: Purely Functional Sets as Characteristic Functions}

To synthesize everything we have learned about functions as first-class citizens, closures, and higher-order functions, we will explore a remarkable computer science paradigm: \textbf{representing sets not as collections of data, but as functions}.\footnote{This classic formulation of purely functional sets as characteristic functions is inspired by and modeled after Prof.~Martin Odersky's landmark course \textit{Functional Programming Principles in Scala} (EPFL).}

In traditional imperative programming, a set is represented by a physical data structure---a hash table, a balanced binary search tree, or a dynamic array. To check if an element belongs to the set, the runtime traverses or hashes memory locations to locate stored items.

In pure functional programming, we can throw away memory buffers entirely. Instead, a set of integers is represented as a \textbf{characteristic function} (or predicate):
\[ \text{FunSet} \triangleq \mathbb{Z} \to \text{Boolean} \]
\textit{Notation Note:} The symbol $\triangleq$ (read as ``is defined to be'') is used throughout formal mathematics and computer science to declare a definition. Unlike a standard equality sign ($=$) which asserts that two existing expressions evaluate to the same value, $A \triangleq B$ means we are defining the term $A$ to be the expression $B$. Here, $\text{FunSet}$ is defined as the type of any function mapping an integer to a boolean.

There is no underlying list or tree. ``Building a set'' means constructing a new predicate function, and ``combining sets'' means composing functions with boolean operators. Membership is computed on demand:
\[ \text{contains}(s, x) \triangleq s(x) \]

\subsection{Categorical Perspective: The Subobject Classifier}

Mathematically, this representation is rooted in \textbf{Topos Theory} and Category Theory. 

In the category $\mathbf{Set}$, any subset $S \subseteq A$ is a \textit{subobject}. A key theorem of category theory states that every subobject $S \subseteq A$ corresponds uniquely to a morphism $\chi_S : A \to \Omega$, where $\Omega = \{\text{True}, \text{False}\}$ is the \textbf{subobject classifier}.

\[
\chi_S(x) = \begin{cases} 
\text{True} & \text{if } x \in S \\ 
\text{False} & \text{if } x \notin S 
\end{cases}
\]

Because the function space $\Omega^A$ (or $A \to \text{Boolean}$) is an exponential object within a Cartesian closed category, the power set $\mathcal{P}(A)$ is isomorphic to the function space $\Omega^A$:
\[ \mathcal{P}(A) \cong \Omega^A \]

This isomorphism tells us that a set and its membership predicate are mathematically identical.

\subsection{Infinite Sets at Zero Memory Cost}

One immediate advantage of characteristic functions over eager data structures is the ability to represent \textbf{infinite sets} in $O(1)$ memory.

For instance, consider the set of all even integers:
\[ E = \{ x \in \mathbb{Z} \mid x \pmod 2 = 0 \} \]
In imperative programming, storing $E$ would require infinite memory. In functional programming, $E$ is simply the closure:
\[ \texttt{evenSet} = \lambda x.\, (x \pmod 2 == 0) \]
Checking membership $\text{contains}(\texttt{evenSet}, 10^9)$ completes instantly without allocating a single byte of memory.

\subsection{Basic Set Operations as Closure Combinators}

Each set operation takes one or two \texttt{FunSet} functions and returns a \textbf{new function} that closes over its inputs. None of these operations loop over elements directly; they wrap the input predicates inside new boolean logic.

\begin{itemize}
    \item \textbf{Singleton Set}: $\text{singletonSet}(e) \triangleq \lambda x.\, (x == e)$
    \item \textbf{Union}: $\text{union}(s, t) \triangleq \lambda x.\, (s(x) \lor t(x))$
    \item \textbf{Intersection}: $\text{intersect}(s, t) \triangleq \lambda x.\, (s(x) \land t(x))$
    \item \textbf{Difference}: $\text{diff}(s, t) \triangleq \lambda x.\, (s(x) \land \neg t(x))$
    \item \textbf{Filter}: $\text{filter}(s, p) \triangleq \lambda x.\, (s(x) \land p(x))$
\end{itemize}

\subsection{Bounded Search: \texttt{forall} and \texttt{exists}}

Because a \texttt{FunSet} is a function, we cannot directly inspect or enumerate its contents. To answer universal or existential questions over a set, we must define a bounded search space $[-K, K]$ (for instance, $K = 1000$).

\paragraph{Universal Quantification (\texttt{forall})}
The expression $\text{forall}(s, p)$ holds true if and only if every element \textit{that is actually in $s$} satisfies predicate $p$:
\[ \forall a \in [-K, K], \quad a \in s \implies p(a) \]
Using material implication ($A \implies B \equiv \neg A \lor B$), the per-element condition is $\neg s(a) \lor p(a)$.

We implement this search using linear tail-recursion scanning from $-K$ up to $K$.

\paragraph{Existential Quantification (\texttt{exists})}
Instead of writing a new loop, we derive $\text{exists}(s, p)$ from \texttt{forall} via De Morgan's Law:
\[ \exists x \in s.\, p(x) \iff \neg (\forall x \in s.\, \neg p(x)) \]

In code:
\[ \text{exists}(s, p) \triangleq \neg\,\text{forall}(s, \lambda x.\, \neg p(x)) \]

\subsection{Transforming Sets: \texttt{map}}

To compute the mapped set $\{ f(x) \mid x \in s \}$, we cannot invert arbitrary functions $f$ to check if $f^{-1}(y) \in s$. Instead, we ask: \textit{does there exist an element $x \in s$ such that $f(x) == y$?}

\[ \text{map}(s, f) \triangleq \lambda y.\, \text{exists}(s, \lambda x.\, (f(x) == y)) \]

This reuses our existential quantifier seamlessly, illustrating how complex set transformations emerge purely by composing higher-order primitives.

\subsection{Implementation in Python and Haskell}

Let us examine the complete code implementation in Python and Haskell.

\lstinputlisting[language=Python, caption={Python: Purely Functional Sets as Characteristic Functions}]{../code/python/chapter4/func_set.py}

\lstinputlisting[language=Haskell, caption={Haskell: Purely Functional Sets as Characteristic Functions}]{../code/haskel/chapter4/FuncSet.hs}
